\textbf{Phần 1. Trắc nghiệm lựa chọn}
\setcounterpageref{ex}{0}

\begin{ex} %[1D9V36]
Mỗi ngày mới bắt đầu từ mấy giờ đêm (rạng sáng)?
\choice
{$1\text{ giờ}$}
{\True $0\text{ giờ}$}
{$2\text{ giờ}$}
{$3\text{ giờ}$}
\loigiai{
Một ngày mới bắt đầu từ 0 giờ đêm (12 giờ đêm).
\textbf{Đáp số:} $0$.\par
\textbf{Đáp án của câu hỏi là B}
}
\end{ex}

\begin{ex} %[1D9V36]
Bố đi làm bằng xe máy lúc 7 giờ sáng và đến công ty lúc 8 giờ sáng. Bố đi trên đường mất bao nhiêu giờ?
\choice
{\True $1$}
{$0$}
{$2$}
{$3$}
\loigiai{
Thời gian đi đường là: $8 - 7 = 1$ (giờ).
\textbf{Đáp số:} $1$.\par
\textbf{Đáp án của câu hỏi là A}
}
\end{ex}

\begin{ex} %[1D9V36]
Buổi trưa, cả nhà em bắt đầu ăn cơm lúc 11 giờ và ăn xong lúc 12 giờ trưa. Bữa cơm trưa kéo dài bao nhiêu giờ?
\choice
{$0$}
{$2$}
{\True $1$}
{$3$}
\loigiai{
Bữa cơm kéo dài: $12 - 11 = 1$ (giờ).
\textbf{Đáp số:} $1$.\par
\textbf{Đáp án của câu hỏi là C}
}
\end{ex}

\begin{ex} %[1D9V36]
Giờ ngủ trưa ở trường bán trú bắt đầu từ 12 giờ trưa đến 2 giờ chiều. Các bạn được ngủ trưa trong bao nhiêu giờ?
\choice
{$1$}
{$3$}
{$4$}
{\True $2$}
\loigiai{
Thời gian ngủ trưa là: $2 - 0 = 2$ (giờ) (từ 12h đến 2h chiều là 2 tiếng).
\textbf{Đáp số:} $2$.\par
\textbf{Đáp án của câu hỏi là D}
}
\end{ex}

\begin{ex} %[1D9V36]
Một tờ lịch tháng ghi: Ngày 5 là Thứ Hai. Hỏi ngày 12 cùng tháng đó là thứ mấy? (Thứ Hai: 2; Thứ Ba: 3).
\choice
{$1$}
{\True $2$}
{$3$}
{$4$}
\loigiai{
Vì $5 + 7 = 12$ nên ngày 12 cũng là Thứ Hai (ghi số 2).
\textbf{Đáp số:} $2$.\par
\textbf{Đáp án của câu hỏi là B}
}
\end{ex}

\begin{ex} %[1D9V36]
Tờ lịch tháng ghi: Ngày 8 là Thứ Sáu. Hỏi ngày 15 cùng tháng đó là thứ mấy? (Thứ Sáu: 6; Thứ Bảy: 7).
\choice
{\True $6$}
{$5$}
{$7$}
{$4$}
\loigiai{
Vì $8 + 7 = 15$ nên ngày 15 cũng là Thứ Sáu (ghi số 6).
\textbf{Đáp số:} $6$.\par
\textbf{Đáp án của câu hỏi là A}
}
\end{ex}

\begin{ex} %[1D9V36]
Nếu ngày 3 là Thứ Bảy thì ngày 10 cùng tháng đó là thứ mấy? (Thứ Bảy: 7; Chủ Nhật: 8).
\choice
{$6$}
{$8$}
{\True $7$}
{$5$}
\loigiai{
Vì $3 + 7 = 10$ nên ngày 10 cũng là Thứ Bảy (ghi số 7).
\textbf{Đáp số:} $7$.\par
\textbf{Đáp án của câu hỏi là C}
}
\end{ex}

\begin{ex} %[1D9V36]
Nếu ngày 17 là Thứ Tư thì ngày 10 cùng tháng đó là thứ mấy? (Thứ Ba: 3; Thứ Tư: 4).
\choice
{$3$}
{$5$}
{$2$}
{\True $4$}
\loigiai{
Vì $17 - 7 = 10$ nên ngày 10 cũng là Thứ Tư (ghi số 4).
\textbf{Đáp số:} $4$.\par
\textbf{Đáp án của câu hỏi là D}
}
\end{ex}

\begin{ex} %[1D9V36]
Đồng hồ reo lúc 6 giờ sáng. Nam thức dậy, vệ sinh cá nhân và ăn sáng mất 1 giờ. Hỏi Nam sẵn sàng đi học lúc mấy giờ?
\choice
{$6\text{ giờ}$}
{\True $7\text{ giờ}$}
{$8\text{ giờ}$}
{$5\text{ giờ}$}
\loigiai{
Nam sẵn sàng đi học lúc: $6 + 1 = 7$ (giờ sáng).
\textbf{Đáp số:} $7$.\par
\textbf{Đáp án của câu hỏi là B}
}
\end{ex}

\begin{ex} %[1D9V36]
Chuyến tàu hỏa khởi hành từ ga lúc 8 giờ tối và chạy suốt đêm đến 6 giờ sáng hôm sau. Hỏi chuyến tàu chạy hết bao nhiêu giờ?
\choice
{\True $10$}
{$0$}
{$20$}
{$9$}
\loigiai{
Từ 8 giờ tối (20h) đến 12h đêm là 4 giờ. Từ 12h đêm đến 6h sáng là 6 giờ. Tổng cộng: $4 + 6 = 10$ (giờ).
\textbf{Đáp số:} $10$.\par
\textbf{Đáp án của câu hỏi là A}
}
\end{ex}

\begin{ex} %[1D9V36]
Bạn Hùng có lịch tập bơi vào Thứ Ba, Thứ Năm và Thứ Bảy hàng tuần. Hỏi trong một tuần Hùng đi tập bơi bao nhiêu buổi?
\choice
{$2$}
{$4$}
{\True $3$}
{$1$}
\loigiai{
Có 3 buổi tập bơi: Thứ Ba, Thứ Năm, Thứ Bảy (tất cả 3 buổi).
\textbf{Đáp số:} $3$.\par
\textbf{Đáp án của câu hỏi là C}
}
\end{ex}

\begin{ex} %[1D9V36]
Tháng có 4 tuần lễ trọn vẹn và thừa 2 ngày. Hỏi tháng đó có bao nhiêu ngày?
\choice
{$20\text{ ngày}$}
{$40\text{ ngày}$}
{$29\text{ ngày}$}
{\True $30\text{ ngày}$}
\loigiai{
4 tuần lễ là: $7 \times 4 = 28$ ngày. Thêm 2 ngày: $28 + 2 = 30$ (ngày).
\textbf{Đáp số:} $30$.\par
\textbf{Đáp án của câu hỏi là D}
}
\end{ex}

\textbf{Phần 2. Trắc nghiệm đúng sai}
\setcounterpageref{ex}{0}

\begin{ex} %[1D9V36]
Xét tính đúng hoặc sai của các khẳng định toán học sau đây:
\choiceTF
{Kết quả của bài toán «Nếu hôm nay là ngày 14 thì 7 ngày trước là ngày mùng mấy» là $6$.}
{\True Kết quả của bài toán «Nếu hôm nay là ngày 18 thì 7 ngày sau là ngày bao nhiêu» là $25$.}
{\True Kết quả của bài toán «Lớp 1A đi tham quan dã ngoại lúc 7 giờ sáng và trở về trường lúc 4 giờ chiều (tức 16 giờ). Chuyến dã ngoại kéo dài bao nhiêu giờ» là $9$.}
{Một ngày có 24 giờ. Em ngủ 9 giờ mỗi ngày. thời gian em thức trong ngày là $14$ giờ.}
\loigiai{
\begin{itemize}
    \item Ý a) Sai: Kết quả đúng là $7$, không phải $6$. Bảy ngày trước là ngày: $14 - 7 = 7$.
\textbf{Đáp số:} $7$.
    \item Ý b) Đúng: Bảy ngày sau là ngày: $18 + 7 = 25$.
\textbf{Đáp số:} $25$.
    \item Ý c) Đúng: Từ 7 giờ sáng đến 4 giờ chiều là: $16 - 7 = 9$ (giờ).
\textbf{Đáp số:} $9$.
    \item Ý d) Sai: Kết quả đúng là $15$, không phải $14$. Thời gian thức là: $24 - 9 = 15$ (giờ).
\textbf{Đáp số:} $15$.
\end{itemize}
\begin{itemize}
    \item a) \textbf{Sai}.
    \item b) \textbf{Đúng}.
    \item c) \textbf{Đúng}.
    \item d) \textbf{Sai}.
\end{itemize}
}
\end{ex}

\begin{ex} %[1D9V36]
Xét tính đúng hoặc sai của các bài toán và phép tính sau đây:
\choiceTF
{\True Kết quả của bài toán «Nếu Thứ Hai là ngày 2 thì Chủ Nhật của tuần đó là ngày mùng mấy» là $8\text{ ngày}$.}
{Kết quả của bài toán «Nếu Chủ Nhật là ngày 9 thì Thứ Hai của tuần đó là ngày mùng mấy» là $2\text{ ngày}$.}
{\True Kết quả của bài toán «Một tiết học kéo dài khoảng 35 phút. Giờ ra chơi kéo dài khoảng 15 phút. Tổng thời gian một tiết học và một giờ ra chơi là bao nhiêu phút» là $50$.}
{Kết quả của bài toán «Lớp 1A đi tham quan dã ngoại lúc 7 giờ sáng và trở về trường lúc 4 giờ chiều (tức 16 giờ). Chuyến dã ngoại kéo dài bao nhiêu giờ» là $8$.}
\loigiai{
\begin{itemize}
    \item Ý a) Đúng: Từ Thứ Hai đến Chủ Nhật có 6 khoảng ngày: $2 + 6 = 8$.
\textbf{Đáp số:} $8$.
    \item Ý b) Sai: Kết quả đúng là $3$, không phải $2$. Thứ Hai đầu tuần đó là ngày: $9 - 6 = 3$.
\textbf{Đáp số:} $3$.
    \item Ý c) Đúng: Tổng thời gian là: $35 + 15 = 50$ (phút).
\textbf{Đáp số:} $50$.
    \item Ý d) Sai: Kết quả đúng là $9$, không phải $8$. Từ 7 giờ sáng đến 4 giờ chiều là: $16 - 7 = 9$ (giờ).
\textbf{Đáp số:} $9$.
\end{itemize}
\begin{itemize}
    \item a) \textbf{Đúng}.
    \item b) \textbf{Sai}.
    \item c) \textbf{Đúng}.
    \item d) \textbf{Sai}.
\end{itemize}
}
\end{ex}

\textbf{Phần 3. Trả lời ngắn}
\setcounterpageref{ex}{0}

\begin{ex} %[1D9V36]
Đồng hồ điện tử hiện số: 08:00. Lúc đó là mấy giờ đúng?
\shortans{8}
\loigiai{
\textbf{Đáp số:} $8$\par
08:00 tương ứng là 8 giờ đúng.
\textbf{Đáp số:} $8$.
}
\end{ex}

\begin{ex} %[1D9V36]
Đồng hồ điện tử hiện số: 17:00. Lúc đó là mấy giờ chiều?
\shortans{5}
\loigiai{
\textbf{Đáp số:} $5$\par
17 giờ tương ứng với 5 giờ chiều (vì $17 - 12 = 5$).
\textbf{Đáp số:} $5$.
}
\end{ex}

\begin{ex} %[1D9V36]
Đồng hồ điện tử hiện số: 20:00. Lúc đó là mấy giờ tối?
\shortans{8}
\loigiai{
\textbf{Đáp số:} $8$\par
20 giờ tương ứng với 8 giờ tối (vì $20 - 12 = 8$).
\textbf{Đáp số:} $8$.
}
\end{ex}

\begin{ex} %[1D9V36]
Tháng Hai năm thường có 28 ngày. Tháng Hai năm thường có bao nhiêu tuần lễ trọn vẹn?
\shortans{4}
\loigiai{
\textbf{Đáp số:} $4$\par
Vì $4 \times 7 = 28$ nên có đúng $4$ tuần lễ trọn vẹn.
\textbf{Đáp số:} $4$.
}
\end{ex}

\begin{ex} %[1D9V36]
Bé được nghỉ tết trong 1 tuần lễ và 3 ngày. Hỏi bé được nghỉ tết tất cả bao nhiêu ngày?
\shortans{10}
\loigiai{
\textbf{Đáp số:} $10$\par
Tổng số ngày nghỉ tết là: $7 + 3 = 10$ (ngày).
\textbf{Đáp số:} $10$.
}
\end{ex}

\begin{ex} %[1D9V36]
Kim giờ quay được 1 vòng quanh mặt đồng hồ (từ số 12 đến số 12) hết bao nhiêu giờ?
\shortans{12}
\loigiai{
\textbf{Đáp số:} $12$\par
Kim giờ quay đúng 1 vòng mặt đồng hồ hết $12$ giờ.
\textbf{Đáp số:} $12$.
}
\end{ex}
